\documentclass[11pt]{article}
\usepackage[T1]{fontenc}
\usepackage[utf8]{inputenc}
\usepackage{amsmath,amssymb,amsthm}
\usepackage{hyperref}

\title{Correction de l'automate et de l'instrumentation}
\author{}
\date{}

\newcommand{\Tr}{\mathcal{T}}
\newcommand{\AP}{\mathsf{AP}}
\newcommand{\eval}{\models}
\newcommand{\Prog}{\mathrm{Prog}}
\newcommand{\nnf}{\mathrm{nnf}}
\newcommand{\res}{\mathrm{res}}
\newcommand{\step}{\mathrm{step}}
\newcommand{\bad}{\mathsf{bad}}
\newcommand{\states}{\mathsf{States}}
\newcommand{\deltaa}{\delta}
\newcommand{\atoms}{\mathsf{Atoms}}

\theoremstyle{definition}
\newtheorem{definition}{Definition}
\newtheorem{lemma}{Lemma}
\newtheorem{theorem}{Theorem}
\newtheorem{corollary}{Corollary}

\begin{document}
\maketitle

\section{Preliminaires}

\subsection{Traces et semantique LTL (fragment $G/X$)}

\begin{definition}[Etats et traces]
Soit $\Sigma$ l'ensemble des etats. Une \emph{trace} est une suite infinie
$\tau = s_0 s_1 s_2 \ldots$ avec $s_i \in \Sigma$. On note $\tau[i] = s_i$.
\end{definition}

\begin{definition}[Atomes]
Soit $\AP$ un ensemble fini d'atomes propositionnels. Une valuation
$\sigma : \AP \to \{\mathsf{true},\mathsf{false}\}$ associe une valeur a chaque atome.
Chaque etat $s \in \Sigma$ definit une valuation $\sigma_s$ des atomes.
\end{definition}

\begin{definition}[Fragment $G/X$]
Les formules LTL sont definies par:
\[
\varphi ::= \mathsf{true} \mid \mathsf{false} \mid p \in \AP \mid \neg\varphi
\mid (\varphi \wedge \varphi) \mid (\varphi \vee \varphi)
\mid X\varphi \mid G\varphi.
\]
Les formules sont en NNF (negations uniquement sur les atomes).
\end{definition}

\begin{definition}[Semantique]
Pour une trace $\tau$ et un indice $i$:
\begin{align*}
\tau,i \eval p &\iff \sigma_{\tau[i]}(p) = \mathsf{true} \\
\tau,i \eval \neg p &\iff \sigma_{\tau[i]}(p) = \mathsf{false} \\
\tau,i \eval \varphi \wedge \psi &\iff \tau,i \eval \varphi \text{ et } \tau,i \eval \psi \\
\tau,i \eval \varphi \vee \psi &\iff \tau,i \eval \varphi \text{ ou } \tau,i \eval \psi \\
\tau,i \eval X\varphi &\iff \tau,i+1 \eval \varphi \\
\tau,i \eval G\varphi &\iff \forall j \ge i,\ \tau,j \eval \varphi.
\end{align*}
On note $\tau \eval \varphi$ pour $\tau,0 \eval \varphi$.
\end{definition}

\subsection{Progression}

\begin{definition}[Progression]
La progression $\Prog(\varphi,\sigma)$ renvoie une formule residuelle telle que:
\[
\tau,i \eval \varphi \iff \tau,i+1 \eval \Prog(\varphi,\sigma_{\tau[i]}).
\]
Sur le fragment $G/X$ (NNF):
\[
\Prog(p,\sigma)=\begin{cases}\mathsf{true} & \sigma(p)=\mathsf{true}\\ \mathsf{false} & \text{sinon}\end{cases}
\quad
\Prog(\neg p,\sigma)=\begin{cases}\mathsf{true} & \sigma(p)=\mathsf{false}\\ \mathsf{false} & \text{sinon}\end{cases}
\]
\[
\Prog(\varphi\wedge\psi,\sigma)=\Prog(\varphi,\sigma)\wedge\Prog(\psi,\sigma),\quad
\Prog(\varphi\vee\psi,\sigma)=\Prog(\varphi,\sigma)\vee\Prog(\psi,\sigma)
\]
\[
\Prog(X\varphi,\sigma)=\varphi,\quad
\Prog(G\varphi,\sigma)=\Prog(\varphi,\sigma)\wedge G\varphi.
\]
\end{definition}

\section{Automate residual}

\begin{definition}[Residus et automate]
Soit $\varphi$ une formule en NNF. On definit l'ensemble fini $\states$ des residus
comme l'ensemble des formules obtenues en appliquant iterativement $\Prog$
a partir de $\varphi$, puis en simplifiant. L'automate residual est:
\[
\mathcal{A}_\varphi = (\states, q_0, \deltaa, \bad)
\]
avec $q_0 = \varphi$, $\bad = \mathsf{false}$, et
$\deltaa(q,\sigma)=\mathrm{Simplify}(\Prog(q,\sigma))$.
\end{definition}

\begin{lemma}[Correctness de la progression]
Pour toute trace $\tau$, tout $i$ et toute formule $\varphi$:
\[
\tau,i \eval \varphi \iff \tau,i+1 \eval \Prog(\varphi,\sigma_{\tau[i]}).
\]
\end{lemma}
\begin{proof}
Par induction structurelle sur $\varphi$, en utilisant la definition de $\Prog$.
\end{proof}

\begin{theorem}[Correction de l'automate]
Pour toute trace $\tau$, la suite d'etats $q_{i+1}=\deltaa(q_i,\sigma_{\tau[i]})$
avec $q_0=\varphi$ verifie:
\[
\tau \eval \varphi \iff \forall i,\ q_i \ne \bad.
\]
\end{theorem}
\begin{proof}[Esquisse]
Par induction sur $i$ et le lemme de progression, on montre que
$\tau,i \eval q_i$ pour tout $i$. Ainsi, si un residu vaut $\mathsf{false}$,
la formule est violee. Reciprocquement, si aucun residu n'est $\mathsf{false}$,
la formule est satisfaite.
\end{proof}

\section{Instrumentation du programme}

\subsection{Langage imperatif jouet}

\begin{definition}[Programme]
Un programme est un systeme de transition $(\Sigma,\rightarrow)$ avec etat initial
$s_0$. Une execution est une trace $\tau$ telle que $s_0 \rightarrow s_1 \rightarrow \cdots$.
\end{definition}

\begin{definition}[Instrumentation moniteur]
On etend l'etat a $\Sigma'=\Sigma \times \states$.
L'etat initial est $(s_0,q_0)$. La transition est:
\[
(s,q)\rightarrow (s',q') \iff s\rightarrow s' \text{ et } q'=\deltaa(q,\sigma_s).
\]
On ajoute l'invariant $q \ne \bad$.
\end{definition}

\begin{theorem}[Correction de l'instrumentation]
Pour toute execution instrumentee $\tau'=(s_0,q_0)(s_1,q_1)\ldots$:
\[
\forall i,\ q_i \ne \bad \iff s_0 s_1 s_2 \ldots \eval \varphi.
\]
Donc, prouver l'invariant $q \ne \bad$ sur $\Sigma'$ garantit la propriete LTL.
\end{theorem}
\begin{proof}[Esquisse]
La projection de $\tau'$ sur $\Sigma$ donne une trace $\tau$ du programme.
La suite $(q_i)$ suit l'automate residual par construction. On conclut par
la correction de l'automate.
\end{proof}

\section{Conclusion}

La preuve se decouple en deux etapes: (1) correction de l'automate residual vis-a-vis
de la formule LTL, (2) correction de l'instrumentation qui synchronise le programme
et l'automate. Ainsi, la satisfaction de l'invariant ``$\bad$ non atteint'' implique
la satisfaction de la specification LTL sur toutes les executions du programme.

\end{document}
